\section{Coefficient two with a nontrivial residual conductor-seven factor}
\label{sec:coefficient-two-seven-strict}

\paragraph{The original coefficient and the exact family.}
Put $K=\mathbb Q(\omega)=\mathbb Q(\sqrt{-3})$.
Let $\Xi_7^{(3)}$ be the four Dirichlet characters modulo seven
whose orders are three or six.  A character of order six is retained
as that full character, including its quadratic part.
Let $\delta$ be any Dirichlet character of two-power conductor,
including the trivial character.  Let $O$ be the original ring of
integers of a finite extension $E/\mathbb Q_2$ containing the values
of $\delta$ and $\xi\in\Xi_7^{(3)}$, with uniformizer $\varpi$ and
residue field $k$.  No restriction is imposed on its absolute ramification.
Write $c_2$ for the cyclotomic character and $\omega_2$ for its
sign character modulo four.  For $e\geq1$, $s\in\mathbb Z$, and a
continuous $\kappa:G_K\to1+\mathfrak m_O$ unramified away two, put
\[
 D=(E/O)(\delta\xi\,\omega_2^{-s}\kappa^{-1}c_2^{-e}).
 \tag{G184.1}
\]
All finite characters in this formula are restricted to $G_K$.
The full strict condition means zero in $H^1(K_v,D)$ at every
$v\mid2$ and zero in $H^1(I_v,D)$ at every other finite place.
The same definition applies to $D[\varpi]$.

\begin{theorem}[AC-G184: original-coefficient strict zero for four seven factors]
\label{AC-G184}
For every choice just specified,
\[
           \mathcal S_2(K,D[\varpi])=0,
           \qquad \mathcal S_2(K,D)=0.                       \tag{G184.2}
\]
These statements also hold for the actual groups on
$G_{K,\Sigma_2}$ in (FC.2), with their unchanged full local and
inertia conditions.  The only class computation used below is a
completed quotient certificate for the fixed degree-six field
\[
 L=K\,\mathbb Q(\zeta_7+\zeta_7^{-1}),\qquad
                     \operatorname{Cl}(L)=0.                 \tag{G184.3}
\]
It is not a complete-unit-basis assertion.

On each actual host page $h\geq1$, all products
$(\delta\xi)|_{G_K}$ with $\delta$ a character modulo $2^{h+2}$
and $\xi\in\Xi_7^{(3)}$ occur and are distinct.  Their number is
\[
                         4\cdot2^{h+1}.                      \tag{G184.4}
\]
Every one has strict zero for every positive Tate weight on the
prescribed original $O_2$.  Thus the number of actual triples at
$1\leq e<T$ is $4\cdot2^{h+1}(T-1)$ for positive integer $T$.
These labels are disjoint from the pure two-power family with
trivial seven factor.  No assertion is made here for the seven
character of order two, or for a nontrivial cubic factor from $K(\sqrt[3]{269})/K$
or an auxiliary $q,d$ factor.  No derived determinant index is asserted zero.
\end{theorem}

\begin{proof}
\proofstep{1. The residual character and its exact splitting field.}
Every two-power root of unity reduces to one in characteristic two.
The same is true of $\omega_2$, the reduction of $c_2$, and the
specified near-one character $\kappa$.  Therefore
\[
                         D[\varpi]\simeq k(\alpha),
                         \qquad\alpha=\overline\xi|_{G_K}.
 \tag{G184.5}
\]
The character $\alpha$ has exact order three: reduction preserves
prime-to-two roots, and kills only the quadratic part of an
order-six value.  Both possible nontrivial cubic reductions of
characters modulo seven have kernel $\{1,-1\}$ in
$(\mathbb Z/7\mathbb Z)^\times$.
Their common splitting field over $\mathbb Q$ is
$F=\mathbb Q(\zeta_7+\zeta_7^{-1})$, a real cyclic cubic field.
Since $K$ is imaginary quadratic, $K\cap F=\mathbb Q$ and
$\operatorname{Gal}(L/K)=C_3$.
The original residue field already contains the values of $\alpha$;
no coefficient extension is made.

Only the residual character is killed over $L$ in this argument.
For an order-six $\xi$, its quadratic factor can remain ramified at
seven over $L$.  We neither discard it nor apply an
unramified-outside-two pro-two theorem to that full coefficient.
The original local argument is carried out over $K$ in Step 4.

\proofstep{2. Identification and completed class-quotient certificate.}
Let $a=\zeta_7+\zeta_7^{-1}$.  The cyclotomic relation gives
\[
                         f(a)=0,\qquad f(Y)=Y^3+Y^2-2Y-1.
\]
The rational-root test proves that $f$ is irreducible.
For $z=a+\sqrt{-3}$ compute the exact resultant
\[
 \begin{split}
 P(X)&=\operatorname{Res}_Y\bigl(f(Y),(X-Y)^2+3\bigr)\\
 &=X^6+2X^5+6X^4+6X^3+35X^2+28X+91.
 \end{split}                                                 \tag{G184.6}
\]
The six conjugates $a_i\pm\sqrt{-3}$ are distinct: all $a_i$ are
real and distinct, whereas the two signs have opposite nonzero
imaginary part.  Since $L/\mathbb Q$ has degree six, $z$ has degree
six and generates exactly $L$.  Thus (G184.6) identifies the input
field before any computer class calculation.

The retained fresh PARI/GP 2.19.1 input constructs this resultant
and runs $B=\operatorname{bnfinit}(P,1)$ at 128 decimal digits.
The exact recorded output is
\[
 \begin{array}{c|c}
 \operatorname{polisirreducible}(P)&1\\
 B.\mathrm{sign}&[0,3]\\
 B.\mathrm{disc}&-64827=-3^3\,7^4\\
 B.\mathrm{no}&1\\
 B.\mathrm{cyc}&[\,]\\
 \operatorname{bnfcertify}(B,1)&1.
 \end{array}                                                 \tag{G184.7}
\]
The preliminary $B.\mathrm{no}$ and $B.\mathrm{cyc}$ are not used
as an unconditional certificate by themselves.
The documented flag-one conclusion is that the actual class group
is a quotient of the computed group.  That group is trivial, so
the actual class group is trivial.  This proves (G184.3), in
particular its two-primary part is zero.

The exact input, output, completed execution receipt and saved
$B$ are retained.  The final run terminated normally with exit 0
and the displayed certificate 1.  An earlier wrapper rejected a
standard maximum-stack warning despite the same successful PARI
result; that attempt is preserved separately.  No unfinished or
conditional computation supplies (G184.3).
The primary software input is the complete PARI manual
\texttt{bnfcertify} clause, including the flag-one scope; the
\texttt{bnfinit} scope and complete \texttt{bnfcertify0} interface
are also retained.  This imports the documented algorithmic
certificate and does not claim a new proof of its entire algorithm
or a full unit certificate.

\proofstep{3. The strict residual group vanishes.}
Put $\Gamma=\operatorname{Gal}(L/K)=C_3$.
Restriction is injective on $H^1(G_K,k(\alpha))$ after passage to
$G_L$: its inflation--restriction kernel is
$H^1(\Gamma,k(\alpha))=0$, since three is invertible in $k$.
The restricted coefficient has trivial action, so its continuous
degree-one cocycles are additive homomorphisms from $G_L$ to $k$.

A strict residual class restricts to an unramified homomorphism at
every finite place of $L$.  At places over two it was already zero
on the entire local group over $K$; elsewhere it kills inertia,
and restriction preserves this.  There are no real places.
The unramified abelian class-field identification therefore gives
an injection
\[
 \mathcal S_2(K,k(\alpha))\lhook\joinrel\longrightarrow
 \operatorname{Hom}_{\mathbb F_2[\Gamma]}
       (\operatorname{Cl}(L)/2,k(\alpha))=0.                  \tag{G184.8}
\]
The equivariance is the usual Artin equivariance; no inverse
character convention can change the conclusion because the
entire group $\operatorname{Cl}(L)/2$ is zero.
In particular this step divides only by three.  It does not
average a two-primary host group or replace the original residue
field by a different coefficient field.

\proofstep{4. Original local invariants, including the order-six case.}
At two, $K_2/\mathbb Q_2$ is unramified quadratic: the polynomial
$X^2+X+1$ is irreducible modulo two.
The character $\xi$ is unramified there, and an arithmetic
Frobenius lift for $K_2$ has value $\xi(4)$.
For all four $\xi$ under consideration its reduction has exact
order three.  Explicitly 3 generates $(\mathbb Z/7)^\times$,
$4=3^4$, and the four exponents of an order-six generator are
$1,2,4,5$, none divisible by three.
Every other factor in (G184.1) reduces to one, even if $\delta$
or $\kappa$ is ramified at two.  Hence the scalar of this lift
minus one is an $O$-unit.  It follows that
\[
                D^{G_{K_2}}=0,\qquad D^{G_K}=0.              \tag{G184.9}
\]

At each place over seven, $K/\mathbb Q$ is unramified (indeed
seven splits), so the local inertia still maps onto
$(\mathbb Z/7)^\times$ in the conductor-seven extension.
Choose an inertia element whose $\overline\xi$-value is a
nontrivial cube root.  All other factors in (G184.1) are
unramified there.  Its scalar minus one is again an $O$-unit,
so
\[
                              D^{I_v}=0\quad(v\mid7).
 \tag{G184.10}
\]
This applies to order-six $\xi$ without killing its quadratic
inertia factor: its nontrivial cubic reduction already gives
the unit difference on the original module over $K$.
At every other finite $v\nmid2$, all factors of (G184.1) are
unramified.  Thus $D^{I_v}=D$, which is divisible by $\varpi$.

Apply the exact sequence
\[
                  0\longrightarrow D[\varpi]
                    \longrightarrow D\xrightarrow{\varpi}D
                    \longrightarrow0.                      \tag{G184.11}
\]
For any of the relevant local groups, the kernel of
$H^1(D[\varpi])\to H^1(D)$ is the image of
$H^0(D)/\varpi H^0(D)$.
Equations (G184.9)--(G184.10) and divisibility at the other
inertia groups make this quotient zero in every case.
Any global class in $\mathcal S_2(K,D)$ killed by $\varpi$
therefore lifts globally to a residual class whose local
restrictions are all zero.  Step 3 makes that lift zero.
Consequently the strict group has no nonzero $\varpi$-torsion.
Every continuous cocycle into the discrete torsion module $D$
has finite image, by compactness of $G_K$, and is therefore
killed by a power of $\varpi$.
A nonzero such class would yield a nonzero class killed by
$\varpi$ after multiplication by a suitable power.
This contradiction proves (G184.2) on the original arbitrary $O$.

\proofstep{5. The literal actual host and all original labels.}
The actual host in (FG.3) contains
$\mathbb Q(\zeta_{2^{h+2}},\zeta_7)$.
The values of a character modulo $2^{h+2}$ have order dividing
$2^h$ for $h\geq1$, and the seven character has order dividing six.
Their product order therefore divides $3\cdot2^h$, precisely the
exponent in (FC.1).  All values already belong to the prescribed
completion $O_2$ of $\mathbb Z[\zeta_{3\cdot2^h}]$.
On that completion the canonical coefficient identification is
\[
 (O_2\otimes_{\mathbb Z_2}\mathbb Q_2/\mathbb Z_2)(\eta)(-e)
                 =(E_2/O_2)(\eta c_2^{-e}),\qquad
                  \eta=(\delta\xi)|_{G_K}.                  \tag{G184.12}
\]
Apply the theorem with this $O_2$, $s=0$ and $\kappa=1$.
Inflation from $G_{K,\Sigma_2}$ to $G_K$ is injective on the
unchanged coefficient, since it is unramified outside that set.
An inflated strict cocycle has zero inertia restriction outside
the set and the identical full local and inertia restrictions
inside it.  Thus the actual (FC.2) strict group injects into
the zero group just proved.

There are $\varphi(2^{h+2})=2^{h+1}$ two-power characters and
four allowed seven characters.  Their rational products are
distinct by their disjoint conductor support.  Their restrictions
to $G_K$ remain distinct: a rational character trivial on $G_K$
is trivial or $\chi_{-3}$, and a character supported only at two
and seven cannot equal the latter.
This proves (G184.4) and its full positive-weight count.
In particular the order-six labels are counted as actual full
characters, not replaced by their cubic reductions.
\end{proof}

\paragraph{Source and scope boundary.}
The class-field equivariance is Sharifi, Proposition 1.3.2, with
the complete containing passage 1.3.1--1.3.5 freshly read; its
class-field foundation is retained.  The original uniformizer
argument is proved explicitly above, with no reflection or
elliptic-unit theorem.  The exact software certificate uses the
documented PARI quotient interface only.
The finite-cover and Artin--Tate dictionary remains the previously
audited SGA 1, Expos\'e V, sections 4,7 and Stacks 03PK scope.
No new archive attribution, motive or Lean execution is claimed.
The theorem gives pointwise strict zeros on an additional actual
family; a same-window numerical reduction still uses the separate
one-union (FL.4) comparison.  The all-prime remainder, mixed
realization $\mathcal R_{\rm cub}$ and H04.4 remain open.
